%% R-BurstHADS --- IEEE two-column manuscript
%% Every number in this file is a cell of RESULTS_PACK.md, cited in comments
%% as [table]. No number appears that is not in the pack.
%%
%% Figures: six, all single column (3.40 in), one graph each, in paper/fig/.
%% Include them with NO width= option; they are already the right size.

\documentclass[conference]{IEEEtran}
\IEEEoverridecommandlockouts

\usepackage{cite}
\usepackage{amsmath,amssymb}
\usepackage{algorithm}
\usepackage{algpseudocode}
\usepackage{graphicx}
\usepackage{booktabs}
\usepackage{url}
\usepackage[hidelinks]{hyperref}

\graphicspath{{fig/}}

\begin{document}

\title{Deadline-Aware Spot Provisioning for Bag-of-Tasks\\
Scheduling under Spot Hibernation}

\author{\IEEEauthorblockN{}\IEEEauthorblockA{}}

\maketitle

\begin{abstract}
Spot virtual machines are cheap but may be hibernated by the provider at any
time, which threatens the deadline of a bag-of-tasks application. Burst-HADS
addresses this by rescuing interrupted tasks onto burstable instances, buying
makespan at a cost premium. We present R-BurstHADS, which adds deadline-aware
preemptive and reactive spot provisioning on top of an otherwise unmodified
Burst-HADS, and evaluate all three schedulers in one discrete-event simulator
frozen before the final runs. The evaluation covers 80 cells --- five
hibernation scenarios, four bag sizes and four deadline factors --- at 30 seeds
each, giving 7{,}200 runs per launch-limit setting. Under the reference
implementation's provider launch limits, R-BurstHADS reduces makespan by
27.9\,\% against HADS and 12.7\,\% against Burst-HADS, while holding its cost
premium over HADS to 12.9\,\% against Burst-HADS's 21.3\,\% --- a reduction of
the premium by 8.4 points, about two fifths of it --- and is 5.4\,\% cheaper
than Burst-HADS. Without launch limits it is 17.8\,\% faster and 14.6\,\%
cheaper than Burst-HADS, dominating it on both axes in 69 of 80 cells. We also
report two negative results in full. Under limits R-BurstHADS is the only
scheduler with more than a single missed task, 109 tasks in 4.3\,\% of runs
against one and none, and all of them disappear when the same runs are repeated
without limits. And our re-implemented baselines do not reproduce the published
hibernation cost result: Burst-HADS costs 18.8\,\% more than HADS in our
simulator against the 1.92\,\% reported. A fidelity audit accompanies the work,
documenting 29 deviations from the published algorithms, of which 10 were
corrected, 13 are deliberate design choices, one is a retained addition and five
remain open.
\end{abstract}

\begin{IEEEkeywords}
cloud scheduling, spot instances, hibernation, burstable instances,
bag-of-tasks, deadline-constrained scheduling, reproducibility
\end{IEEEkeywords}

%% ===================================================================
\section{Introduction}
%% ===================================================================

Infrastructure-as-a-service providers sell unused capacity at a steep discount
on the understanding that they may reclaim it. On Amazon EC2 a spot instance can
be \emph{hibernated}: its memory and context are written to storage, execution
stops, and the instance may or may not return. For a bag-of-tasks application
with a deadline this is the central difficulty. Spot capacity is where the
savings are, and it is also the capacity that disappears.

HADS~\cite{teylo2021} schedules such applications across spot and on-demand
markets, checkpointing tasks and migrating them when an instance is reclaimed.
Burst-HADS~\cite{teylo2023} extends it with burstable instances, which are kept
available and can absorb a rescued task immediately at full speed by spending
accumulated CPU credits, and with a work-stealing procedure that redistributes
queued work when a machine becomes idle. The mechanism buys makespan, and it is
not free: in our measurements under provider launch limits Burst-HADS finishes
19.1\,\% sooner than HADS and pays 21.3\,\% more for it.%
\footnote{Table~\ref{tab:limits-on}, all-cells row.}

That trade is the starting point of this work. Burst-HADS reacts to an
interruption by finding somewhere to put the displaced task among the machines
it already has; when the burstable pool is occupied, the only remaining
destination is freshly rented on-demand capacity, the most expensive thing in
the catalogue. A scheduler that instead anticipates interruptions, and that
replaces lost capacity with more spot capacity rather than renting, should be
able to keep the makespan benefit while giving back much of the premium.

\subsection{Contributions}

\begin{itemize}
\item \textbf{R-BurstHADS}, a scheduler that adds two provisioning tiers around
an unmodified Burst-HADS: a preemptive tier that buys replacement spot capacity
for machines judged likely to be reclaimed, and a reactive tier that sizes new
capacity from the remaining deadline slack when the published rescue attempts
are exhausted (Section~\ref{sec:design}).

\item \textbf{A controlled comparison} of the three schedulers over 80 cells
$\times$ 30 seeds under two account regimes --- with and without the reference
implementation's launch limits --- on the same 7{,}200 units of work
(Sections~\ref{sec:limits-on} and~\ref{sec:limits-off}).

\item \textbf{A fidelity audit.} We validate our re-implementations against the
published results, report where they fail to reproduce them, rule out three
candidate explanations by measurement, and publish a 29-row deviation register
(Section~\ref{sec:audit}). Design decisions with more than one defensible answer
were fixed by a rule committed to version control before the runs that decided
them.

\item \textbf{Negative results reported in full}, including the one regime in
which our scheduler is the worst of the three
(Sections~\ref{sec:misses} and~\ref{sec:limitations}).
\end{itemize}

\subsection{What this paper is not}

It is not a reproduction of Teylo \emph{et al.}~\cite{teylo2023}. Our HADS
implementation lands within one to two percent of the published makespans on
three of four workloads without hibernation, but our Burst-HADS does not
reproduce either its published makespan reduction or, more importantly, its
published hibernation cost behaviour. Section~\ref{sec:audit} reports this in
detail, and every comparison in this paper is a comparison against our
re-implementations as documented rather than against the published systems.

%% ===================================================================
\section{Background}
%% ===================================================================

\subsection{Spot hibernation and the two baselines}

HADS partitions a bag of tasks across a pool of spot and on-demand machines with
a single-pass greedy primary scheduler, maintains a checkpoint for each running
task, and on hibernation migrates the affected tasks by walking a fixed ladder:
another machine already in the fleet first, then an on-demand machine already in
the fleet, and finally a newly launched on-demand machine.

Burst-HADS replaces the primary scheduler with an iterated local search that
minimises a weighted combination of cost and makespan, and inserts burstable
instances into the recovery ladder. A burstable machine holds exactly one task
at a time. Rescued onto one, a task runs in \emph{burst mode} at full speed and
spends CPU credits; picked up by work stealing, it runs in \emph{baseline mode}
at a fraction of full speed and earns them.

\subsection{Burstable economics depend on the catalogue}

Whether rescuing onto a burstable is a saving is a property of the price list,
not of the algorithm. In the catalogue used by~\cite{teylo2023} a burst-mode
task pays $0.83\times$ what the same task pays running alone on a newly launched
on-demand machine, which is the alternative the recovery ladder would otherwise
take. In a catalogue built from current Amazon EC2 prices the same comparison is
$1.15\times$: the burstable instance is \emph{dearer} per unit of work than the
on-demand machine it is meant to displace.%
\footnote{Table~\ref{tab:catalogue}; full price table in the results pack [T15].}
We report results on both catalogues and treat this as a finding rather than a
parameter choice.

%% ===================================================================
\section{R-BurstHADS}
\label{sec:design}
%% ===================================================================

R-BurstHADS inherits Burst-HADS in full --- the same primary scheduler, the same
Algorithm~4 recovery attempts, the same work stealing --- and wraps two
provisioning tiers around them. Nothing in the inherited path is modified, so a
correction to the baseline's fidelity is inherited automatically. Notation is
given in Table~\ref{tab:symbols}.

\begin{table}[!t]
\caption{Notation}
\label{tab:symbols}
\centering
\small
\begin{tabular}{@{}ll@{}}
\toprule
symbol & meaning \\
\midrule
$D$ & application deadline \\
$D_{\mathrm{spot}}$ & $D$ minus the worst-case migration reserve \\
$t$ & current simulation time \\
$s_v$ & per-core speed of machine $v$ \\
$k_v$ & usable execution slots (vCPUs) of $v$ \\
$\theta_v$ & whole-machine throughput, $s_v k_v$ \\
$r_v$ & price of $v$ per unit time \\
$\lambda_v$ & hibernation rate of $v$ \\
$n_v$ & tasks the primary schedule placed on $v$ \\
$\bar{e}$ & mean task execution time \\
$\omega$ & instance start-up latency \\
$\Theta(A)$ & aggregate throughput of the active fleet $A$ \\
$W$ & remaining unfinished work \\
$P$ & set of machines provisioned by R-BurstHADS \\
\bottomrule
\end{tabular}
\end{table}

\subsection{Tier 0: preemptive replacement}

After the primary schedule is built, each spot machine is examined once. A
replacement is bought when the expected cost of reacting to its loss exceeds the
cost of anticipating it:
\begin{equation}
\bigl(1 - e^{-\lambda_v D}\bigr)\, C_{\mathrm{react}}(v) \;>\; C_{\mathrm{proact}},
\label{eq:thm1}
\end{equation}
where the reactive cost is the extra paid to finish $v$'s work on rented
on-demand capacity instead of spot,
\begin{equation}
C_{\mathrm{react}}(v) = n_v\,\bar{e}\left(\frac{r_{\mathrm{od}}}{\theta_{\mathrm{od}}}
- \frac{r_{s}}{\theta_{s}}\right),
\label{eq:creact}
\end{equation}
and the proactive cost is the replacement's idle boot period,
$C_{\mathrm{proact}} = \omega\, r_{s} + \epsilon$. Both rate terms are divided by
\emph{whole-machine} throughput $\theta = s k$, not by per-core speed; dividing
by speed alone prices a four-core machine as though it retired one core's worth
of work.

The replacement's instance type is resolved from the scheduler's own pool, never
from a private constant, so R-BurstHADS can only obtain a type the baselines'
planners could also have selected. Among types that are more likely than not to
survive to the deadline, it takes the best throughput per unit price, breaking
exact ties toward the larger machine.

\subsection{Tier 3: deadline-sized reactive provisioning}

When an interruption occurs and the inherited attempts are exhausted, the
shortfall is sized from what is left of the deadline. With $T = D - t - \omega$
the usable time after a new machine boots, the number of additional machines is
\begin{equation}
n_{\mathrm{extra}} = \left\lceil \frac{W/T - \Theta(A)}{\theta} \right\rceil .
\label{eq:thm2}
\end{equation}
Spot capacity is bought when the remaining slack comfortably exceeds the boot
latency and the task can still finish in time on the new machine; otherwise a
burstable instance is taken, which is usable immediately.

\begin{algorithm}[!t]
\caption{R-BurstHADS: destination selection on hibernation}
\label{alg:select}
\begin{algorithmic}[1]
\Require task $\tau$, time $t$, deadline $D$
\Statex \textbf{Tier 0 --- preemptively provisioned capacity}
\ForAll{$v \in P$ alive}
  \If{\Call{CheckMigration}{$\tau, v, t, D$}} \Return $v$ \EndIf
\EndFor
\Statex \textbf{Tiers 1--2 --- Burst-HADS, unmodified}
\State $v \gets$ \Call{PaperMigration}{$\tau, t$}
  \Comment{idle burstable, then any active machine}
\If{$v \neq \bot$} \Return $v$ \EndIf
\Statex \textbf{Tier 3 --- reactive provisioning}
\State $v \gets$ \Call{ProvisionOneMore}{$\tau, t$}
  \Comment{Eq.~\eqref{eq:thm2}}
\If{$v \neq \bot$} \Return $v$ \EndIf
\Statex \textbf{Tier 4 --- Burst-HADS, unmodified}
\State \Return \Call{OnDemandFallback}{$\tau, t$}
\end{algorithmic}
\end{algorithm}

\begin{algorithm}[!t]
\caption{\textsc{ProvisionOneMore}}
\label{alg:provision}
\begin{algorithmic}[1]
\Require task $\tau$, time $t$
\State $\mathit{slack} \gets D - t$
\If{$\mathit{slack} > 2\omega$}
  \State $f \gets t + \omega + \tau.\mathit{remaining}/s_{\mathrm{tpl}}$
  \If{$f \le D$ \textbf{and} $D - f > \tau.e/s_{\mathrm{tpl}}$}
    \State $v \gets$ \Call{LaunchSpot}{ready at $t+\omega$};
           $P \gets P \cup \{v\}$
    \State \Return $v$
  \EndIf
\EndIf
\If{$t + \tau.\mathit{remaining}/s_{b} \le D$}
  \State $v \gets$ \Call{LaunchBurstable}{ready at $t$};
         $P \gets P \cup \{v\}$
  \State \Return $v$
\EndIf
\State \Return $\bot$
\end{algorithmic}
\end{algorithm}

The guard on line~4 of Algorithm~\ref{alg:provision} is the spare-time rule
of~\cite[\S3.4]{teylo2023}: a spot destination must leave more slack than the
longest task placed on it, so that a second interruption can still be absorbed.
Tier~0 obtains the same rule by calling the inherited \textsc{CheckMigration};
tier~3 states it explicitly because the machine has not yet booted and the
inherited estimator has no notion of boot latency.

%% ===================================================================
\section{Method}
\label{sec:method}
%% ===================================================================

\subsection{Simulator and freeze}

All three schedulers run in one discrete-event simulator. The simulator was
frozen before the final runs at a tagged commit, and every result in this paper
comes from that code; defects found after the freeze are recorded in the
deviation register rather than fixed. Every result file is committed, and the
results pack that generates the tables below records each file's hash and the
commit it came from.

\subsection{Catalogues}

Two instance catalogues are used: a \emph{sweep} catalogue at current EC2 prices
for the main comparison, and a \emph{validation} catalogue reproducing
Table~3 of~\cite{teylo2023} for the fidelity audit. Each has three spot types,
three on-demand types and one burstable type. Key economics are in
Table~\ref{tab:catalogue}.

\begin{table}[!t]
\caption{Price per unit of work. ``1 task'' is what a task pays running alone
on the machine, which is the comparison the recovery ladder makes.}
\label{tab:catalogue}
\centering
\footnotesize
\setlength{\tabcolsep}{3.6pt}
\begin{tabular}{@{}llrrr@{}}
\toprule
catalogue & type (market) & \$/h & \$/unit full & \$/unit, 1 task \\
\midrule
sweep & c5.large (spot) & 0.0306 & 0.00765 & 0.01530 \\
sweep & c5.xlarge (spot) & 0.0612 & 0.00765 & 0.03060 \\
sweep & m5.xlarge (spot) & 0.0700 & 0.01029 & 0.04118 \\
sweep & c5.large (on-dem.) & 0.0850 & 0.02125 & 0.04250 \\
sweep & t3.large (burst.) & 0.0832 & 0.04894 & 0.04894 \\
\midrule
valid. & c4.large (spot) & 0.0366 & 0.01830 & 0.03660 \\
valid. & c4.large (on-dem.) & 0.1000 & 0.05000 & 0.10000 \\
valid. & t3.large (burst.) & 0.0832 & 0.08320 & 0.08320 \\
\bottomrule
\end{tabular}
\end{table}

\subsection{Experiment grid}

The grid crosses the five hibernation scenarios of~\cite[Table~9]{teylo2023} ---
sc1 $(k_h,k_r)=(1,0)$, sc2 $(5,0)$, sc3 $(1,5)$, sc4 $(5,5)$, sc5 $(3,2.5)$,
with per-machine rates $\lambda_h = k_h/D$ and $\lambda_r = k_r/D$ --- with bag
sizes $n \in \{50,100,200,300\}$ and deadline factors
$DF \in \{0.25, 0.5, 1.0, 2.0\}$, at 30 seeds per cell. That is 80 cells and
7{,}200 runs per limit setting. The deadline is
$D = DF \cdot \kappa \cdot M^{*}(n)$ with a floor of 339.7\,s; the floor binds at
$DF = 0.25$ for $n = 100$.

Deadline factor is swept rather than fixed because both baselines only escalate
when the deadline is threatened, so a single deadline setting measures the slack
granted as much as the scheduling policy. Reporting across $DF$ makes the
constant in $D$ non-load-bearing.

\subsection{Launch limits}

The reference implementation limits an account to five launches per instance
type per market and twenty per market, with burstable instances counted against
the on-demand allowance. An uncapped pool lets any deadline be bought and
flattens deadline satisfaction into a metric that cannot discriminate, so we
adopt the limits and additionally report the same 7{,}200 units without them.
The starting pool holds three spot machines per type, below the limit of five,
so the spot limit binds only on capacity R-BurstHADS launches itself. Both
regimes are reported throughout, leading with limits on.

\subsection{Metrics}

Cost is billed per second from launch to shutdown or hibernation,
following~\cite[\S3.1]{teylo2023}: a machine's meter opens when a scheduler
launches it, whether or not it receives work. A cell is a
$(\text{scenario}, n, DF)$ triple; a reported percentage is the ratio of cell
means, averaged over cells. Confidence intervals for R-BurstHADS against
Burst-HADS are paired per seed at 95\,\%. Deadline satisfaction is reported as
two measures --- missed tasks, and the share of runs with any miss --- because
they order the schedulers differently. Runs whose primary scheduler finds no
placement within $D$ and the limits are reported as \emph{infeasible} and
excluded from cross-scheduler means, with their cause.

\subsection{Decision discipline}

Every change to a baseline was first measured as a counterfactual patch and
adopted only when the adopted code reproduced the measured variant row for row:
1{,}440 of 1{,}440 validation rows and 2{,}400 of 2{,}400 sweep rows. Where two
principled options existed, the choice rule was committed to version control
before the runs that would decide it.

%% ===================================================================
\section{Fidelity audit}
\label{sec:audit}
%% ===================================================================

\subsection{What reproduces}

Without hibernation, our HADS lands close to the published makespans:
2330\,s against 2290\,s on J60, 2350 against 2295 on J80, 2356 against 2332 on
J100, and 2365 against 2580 on ED200.

\subsection{What does not}

Our Burst-HADS does not. Its makespan reduction against HADS without
hibernation is $-70.9/-59.8/-50.7/-0.1\,\%$ on the four workloads against the
published $-44.4/-42.1/-28.8/-11.8\,\%$: far stronger than published on the
smaller jobs and flat on the largest. Its J60 cost, however, is within two
percent of the published figure.

Under hibernation neither published aggregate is reproduced
(Table~\ref{tab:validation}). The headline discrepancy is cost: Burst-HADS
costs 18.8\,\% more than HADS in our simulator against the published 1.92\,\%.
The two regimes agree to within 0.14 points, so this is not an artefact of the
launch limits.

\begin{table}[!t]
\caption{Validation against~\cite[Table~9]{teylo2023}, 20 cells, same
definition for both. Premium is a scheduler's cost under hibernation against
its own cost without.}
\label{tab:validation}
\centering
\small
\begin{tabular}{@{}lrrr@{}}
\toprule
quantity & ours, limits off & limits on & \cite{teylo2023} \\
\midrule
Burst-HADS cost vs HADS & $+18.80\,\%$ & $+18.66\,\%$ & $+1.92\,\%$ \\
Makespan reduction & $22.18\,\%$ & $22.18\,\%$ & $25.87\,\%$ \\
\quad J60 only & $35.43\,\%$ & $35.43\,\%$ & $40.10\,\%$ \\
\quad ED200 only & $5.28\,\%$ & $5.25\,\%$ & $10.24\,\%$ \\
HADS premium & $+105\,\%$ & $+105\,\%$ & $+95\,\%$ \\
Burst-HADS premium & $+40\,\%$ & $+40\,\%$ & $+25\,\%$ \\
\bottomrule
\end{tabular}
\end{table}

\subsection{Where the gap is, and what it is not}

The gap is not shared between the schedulers. Hibernation raises HADS's cost by
105\,\% against a published 95\,\%, which is close; it raises Burst-HADS's by
40\,\% against a published 25\,\%. Burst-HADS's premium is the discrepancy.

Three candidate explanations were tested and rejected by measurement. The
\emph{aggregation} is not at fault: alternative weightings of our runs give
$+36.6$ to $+42.4\,\%$ against our headline $+38.7\,\%$ at the time of the test,
all far from $+1.9\,\%$. The \emph{billing rule} is not at fault: no rule closes
the gap, with 900\,s quantisation giving $+30.1\,\%$, billing from launch
$+40.4\,\%$, the published rule $+42.5\,\%$, and an earlier rule $+22.9\,\%$ at
the price of wrecking the no-hibernation validation. \emph{Storage charges
during hibernation} move the figure from $+38.7$ to $+37.5\,\%$.

What remains is exposure. Our runs suffer between 1.3 and 3.6 times the
hibernations per run that the published table implies, because our schedulers
deploy most of the spot pool and hibernation is applied per machine. The pool
size on which those counts depend is not stated in~\cite{teylo2023}, and no
value we tried reproduces its Table~7 and Table~9 together. We record this as a
finding about the published description rather than as a parameter to tune.

\subsection{The audit moved the numbers}

Four baseline corrections adopted during the audit moved the validation toward
the published results: the cost change from $+38.7$ to $+18.8\,\%$, the makespan
reduction from 14.7 to 22.2\,\%, and Burst-HADS's premium from $+94$ to
$+40\,\%$. They did not close the gap.

\subsection{Deviation register}

Twenty-nine deviations from the published algorithms are recorded with their
measured effect: 10 corrected, 13 deliberate design choices, one retained
addition and five open, one of which is unmeasured. The register is reproduced
in Appendix~A.

One entry deserves the main text. Our Burst-HADS keeps an improvement guard in
its proactive burstable allocation which the published description does not
have, and the guard favours the comparison reported here: with it, R-BurstHADS
is 5.3\,\% cheaper than Burst-HADS with 34 of 55 cells dominated; without it,
5.8\,\% dearer with 14 of 55. It was kept by the pre-registered rule, because
removing it turned 85 clean Burst-HADS runs and 92 clean R-BurstHADS runs into
runs with a missed task, against none and three in the other direction.

%% ===================================================================
\section{Results under provider launch limits}
\label{sec:limits-on}
%% ===================================================================

\subsection{Headline}

\begin{figure}[!t]
\centering
\includegraphics{fig1_headline.pdf}
\caption{Makespan and cost against HADS under launch limits, over the 75 cells
feasible for every scheduler in every seed. R-BurstHADS roughly halves
Burst-HADS's cost premium while finishing sooner. Over all 80 cells with
available seeds the makespan figures are 26.4\,\% against HADS and 11.9\,\%
against Burst-HADS.}
\label{fig:headline}
\end{figure}

Over the 75 cells in which every scheduler is feasible in every seed
(Table~\ref{tab:limits-on}, Fig.~\ref{fig:headline}), Burst-HADS reduces
makespan by 19.1\,\% against HADS at 21.3\,\% more cost. R-BurstHADS reduces it
by 27.9\,\% at 12.9\,\% more cost: it cuts the premium by 8.4 points, about two
fifths of it, and is simultaneously 12.7\,\% faster and 5.4\,\% cheaper than
Burst-HADS, dominating it on both axes in 44 of the 75 cells. Paired per-seed
intervals put it significantly faster than Burst-HADS in 45 cells and slower in
9, significantly cheaper in 29 and dearer in 7.

\begin{table}[!t]
\caption{Results under launch limits. B and R are Burst-HADS and R-BurstHADS;
H is HADS. Percentages are ratios of cell means averaged over cells.}
\label{tab:limits-on}
\centering
\footnotesize
\setlength{\tabcolsep}{3.2pt}
\begin{tabular}{@{}lrrrrrrr@{}}
\toprule
group & cells & B mk & B \$ & R mk & R \$ & R mk & R \$ \\
 & & vs H & vs H & vs H & vs H & vs B & vs B \\
\midrule
all & 75 & $-19.1$ & $+21.3$ & $-27.9$ & $+12.9$ & $-12.7$ & $-5.4$ \\
\midrule
$DF{=}0.25$ & 15 & $-7.3$ & $+9.7$ & $-8.9$ & $+8.4$ & $-1.8$ & $-1.3$ \\
$DF{=}0.5$ & 20 & $-6.2$ & $+18.6$ & $-7.4$ & $+18.6$ & $-1.4$ & $-0.0$ \\
$DF{=}1.0$ & 20 & $-13.1$ & $+24.2$ & $-27.0$ & $+16.6$ & $-15.9$ & $-4.3$ \\
$DF{=}2.0$ & 20 & $-46.6$ & $+30.0$ & $-63.6$ & $+7.1$ & $-28.9$ & $-14.9$ \\
\midrule
sc1 & 15 & $-27.1$ & $+36.3$ & $-39.2$ & $+19.8$ & $-18.9$ & $-9.8$ \\
sc2 & 15 & $-13.0$ & $+20.2$ & $-18.2$ & $+18.8$ & $-7.7$ & $-1.3$ \\
sc3 & 15 & $-25.5$ & $+24.8$ & $-37.6$ & $+10.4$ & $-17.4$ & $-9.5$ \\
sc4 & 15 & $-12.7$ & $+12.6$ & $-18.8$ & $+10.2$ & $-7.5$ & $-1.0$ \\
sc5 & 15 & $-16.9$ & $+12.8$ & $-25.8$ & $+5.5$ & $-12.0$ & $-5.3$ \\
\midrule
$n{=}50$ & 20 & $-11.7$ & $+32.4$ & $-28.7$ & $+11.5$ & $-19.2$ & $-12.3$ \\
$n{=}100$ & 15 & $-18.4$ & $+28.1$ & $-31.2$ & $+16.5$ & $-18.9$ & $-8.6$ \\
$n{=}200$ & 20 & $-21.8$ & $+14.5$ & $-25.6$ & $+11.4$ & $-7.7$ & $-2.6$ \\
$n{=}300$ & 20 & $-24.2$ & $+12.1$ & $-27.1$ & $+13.2$ & $-6.5$ & $+1.1$ \\
\midrule
all 80 & 80 & $-18.0$ & $+20.0$ & $-26.4$ & $+12.2$ & $-11.9$ & $-5.0$ \\
\bottomrule
\end{tabular}
\end{table}

\subsection{The advantage needs deadline slack}

\begin{figure}[!t]
\centering
\includegraphics{fig2_cost_vs_df.pdf}
\caption{Cost against HADS as a function of deadline slack, under launch
limits. The two schedulers are indistinguishable at $DF{=}0.5$, both at
$+18.6\,\%$, and separate as slack grows. R-BurstHADS is never cheaper than
HADS in this regime.}
\label{fig:cost-df}
\end{figure}

\begin{figure}[!t]
\centering
\includegraphics{fig3_makespan_vs_df.pdf}
\caption{Makespan against HADS as a function of deadline slack, under launch
limits. Both baselines finish close to the deadline wherever they can buy their
way there, so the gap widens with slack.}
\label{fig:mk-df}
\end{figure}

Figures~\ref{fig:cost-df} and~\ref{fig:mk-df} show that the advantage is
concentrated where the deadline is loose. Against Burst-HADS, R-BurstHADS is
28.9\,\% faster and 14.9\,\% cheaper at $DF = 2.0$, dominating 17 of 20 cells;
15.9\,\% faster and 4.3\,\% cheaper at $DF = 1.0$; and essentially tied at
$DF = 0.5$, at 1.4\,\% faster and equal on cost, with 7 cells where it is
significantly slower and 6 where it is significantly dearer. At $DF = 0.25$ the
difference is 1.8\,\% and 1.3\,\%.

The reason is visible in HADS's own behaviour: it finishes at 93\,\% of the
deadline on average and 99\,\% at $DF = 0.25$. Both baselines escalate only when
the deadline is threatened, so when it is tight they spend whatever is needed
and the three schedulers converge; when it is loose they trade makespan for
cost and R-BurstHADS, whose provisioning is triggered by interruption risk
rather than by deadline pressure alone, does not.

\subsection{By scenario and by bag size}

\begin{figure}[!t]
\centering
\includegraphics{fig4_scenarios.pdf}
\caption{Cost against HADS by hibernation scenario, under launch limits.
sc1 is where R-BurstHADS costs the most relative to HADS and also where it beats
Burst-HADS by the widest margin.}
\label{fig:scenarios}
\end{figure}

Against Burst-HADS the margin is widest in the low-interruption scenarios ---
18.9\,\% faster and 9.8\,\% cheaper in sc1, 17.4\,\% and 9.5\,\% in sc3 --- and
narrowest in the high-interruption ones, 7.5\,\% and 1.0\,\% in sc4 and 7.7\,\%
and 1.3\,\% in sc2, each with three cells where R-BurstHADS is significantly
dearer. Fig.~\ref{fig:scenarios} shows the complementary view against HADS: sc1
is both where R-BurstHADS pays the largest premium over HADS, 19.8\,\%, and
where it beats Burst-HADS by the most.

\begin{figure}[!t]
\centering
\includegraphics{fig5_cost_vs_n.pdf}
\caption{Cost against HADS by bag size, under launch limits. The cost advantage
over Burst-HADS narrows as the bag grows and reverses at $n{=}300$, where
R-BurstHADS is 1.1\,\% dearer. Without limits there is no reversal.}
\label{fig:cost-n}
\end{figure}

The advantage also narrows with scale (Fig.~\ref{fig:cost-n}). At $n = 50$
R-BurstHADS is 19.2\,\% faster and 12.3\,\% cheaper than Burst-HADS, dominating
17 of 20 cells. At $n = 300$ it is 6.5\,\% faster and 1.1\,\% \emph{dearer},
dominating only 5 of 20. The reversal is weak --- three cells significantly
cheaper against two significantly dearer --- and it does not appear without
launch limits, where the same comparison is 12.8\,\% cheaper. We read it as the
launch ceiling binding hardest where the work is largest.

\subsection{Deadline misses}
\label{sec:misses}

\begin{figure}[!t]
\centering
\includegraphics{fig6_missed_tasks.pdf}
\caption{Missed tasks with and without launch limits. R-BurstHADS is the only
scheduler with more than a single missed task under limits, and every miss
disappears when the same runs are repeated without them.}
\label{fig:misses}
\end{figure}

This is where our scheduler is the worst of the three
(Table~\ref{tab:misses}, Fig.~\ref{fig:misses}). Under launch limits HADS misses
no task, Burst-HADS misses one task in one run, and R-BurstHADS misses 109 tasks
in 104 of 2{,}400 runs --- 4.3\,\% of runs, and 0.028\,\% of all tasks. Eighty of
the 109 are at $DF = 0.5$. The worst cells are sc4 and sc5 at $n = 300$,
$DF = 0.5$, with 15 missed tasks each, half a task per run.

\begin{table}[!t]
\caption{Deadline misses, both measures, feasible runs only.}
\label{tab:misses}
\centering
\footnotesize
\setlength{\tabcolsep}{4pt}
\begin{tabular}{@{}llrrrr@{}}
\toprule
limits & scheduler & runs & runs w/ miss & share & missed tasks \\
\midrule
on & HADS & 2390 & 0 & 0.0\,\% & 0 \\
on & Burst-HADS & 2400 & 1 & 0.0\,\% & 1 \\
on & R-BurstHADS & 2400 & 104 & 4.3\,\% & 109 \\
off & HADS & 2400 & 0 & 0.0\,\% & 0 \\
off & Burst-HADS & 2400 & 0 & 0.0\,\% & 0 \\
off & R-BurstHADS & 2400 & 0 & 0.0\,\% & 0 \\
\bottomrule
\end{tabular}
\end{table}

The evidence that the launch ceiling is responsible is the matched
counterfactual, not the limit flag. Repeating the same 7{,}200 units without
limits removes every missed task. It is true that all 104 missing runs had
reached a launch limit, most often the c5.xlarge spot limit in 103 of them; but
so had 93.5\,\% of R-BurstHADS's runs \emph{without} a miss, so reaching a limit
does not by itself identify the failures.

\subsection{Infeasible runs}

Ten runs of 7{,}200 have no feasible primary schedule, all of them HADS, at
$n = 100$ and $DF = 0.25$ where the deadline sits on its 339.7\,s floor, in
seeds 1 and 24 of every scenario. Burst-HADS and R-BurstHADS are feasible in all
2{,}400 runs each. Because a cross-scheduler cell mean is undefined when one
scheduler has no run, those five cells are excluded from
Table~\ref{tab:limits-on}'s main rows and reported separately; the last row
gives the all-80-cell figures over available seeds. Without limits no run is
infeasible.

%% ===================================================================
\section{Results without launch limits}
\label{sec:limits-off}
%% ===================================================================

Repeating the same 7{,}200 units with the account limits removed isolates how
much of R-BurstHADS's behaviour depends on being allowed to launch
(Table~\ref{tab:limits-off}). All 80 cells are feasible and no scheduler misses
a deadline.

\begin{table}[!t]
\caption{Results without launch limits, 80 cells.}
\label{tab:limits-off}
\centering
\footnotesize
\setlength{\tabcolsep}{3.6pt}
\begin{tabular}{@{}lrrrrrr@{}}
\toprule
group & B mk & B \$ & R mk & R \$ & R mk & R \$ \\
 & vs H & vs H & vs H & vs H & vs B & vs B \\
\midrule
all & $-18.1$ & $+21.9$ & $-31.0$ & $+2.6$ & $-17.8$ & $-14.6$ \\
$DF{=}0.25$ & $-6.0$ & $+12.6$ & $-8.2$ & $+8.7$ & $-2.4$ & $-3.4$ \\
$DF{=}0.5$ & $-6.4$ & $+20.5$ & $-13.2$ & $-2.1$ & $-7.2$ & $-18.2$ \\
$DF{=}1.0$ & $-13.1$ & $+24.6$ & $-34.2$ & $+2.3$ & $-24.0$ & $-16.8$ \\
$DF{=}2.0$ & $-46.6$ & $+30.0$ & $-68.5$ & $+1.4$ & $-37.7$ & $-19.9$ \\
\bottomrule
\end{tabular}
\end{table}

R-BurstHADS is then 17.8\,\% faster and 14.6\,\% cheaper than Burst-HADS,
dominating it in 69 of 80 cells, significantly faster in 70 and slower in none,
significantly cheaper in 59 and dearer in 4. Against HADS it is 31.0\,\% faster
at a 2.6\,\% premium.

The contrast identifies what the advantage runs on. Burst-HADS barely notices
the regime change, moving from $-19.1\,\%/+21.3\,\%$ to
$-18.1\,\%/+21.9\,\%$ against HADS, because it never launches enough to reach a
limit. R-BurstHADS changes a great deal: its cost lead over Burst-HADS grows
from 5.4\,\% to 14.6\,\%, its premium over HADS falls from 12.9\,\% to 2.6\,\%,
and its deadline misses vanish. At $DF = 0.5$, where the two are tied under
limits, removing them makes R-BurstHADS 7.2\,\% faster and 18.2\,\% cheaper in
19 of 20 cells.

Headroom, in other words, is the resource this scheduler trades on. Where the
provider grants it the method is strong; where it does not, most of the cost
advantage and all of the reliability advantage go away.

%% ===================================================================
\section{Threats to validity and limitations}
\label{sec:limitations}
%% ===================================================================

\textbf{This is not a reproduction.} The published hibernation cost result is
not reproduced, $+18.8\,\%$ against $+1.92\,\%$, and neither are the published
Burst-HADS makespans without hibernation. Every comparison here is against our
re-implementations as documented in the deviation register, and the residual
gap is concentrated in Burst-HADS's hibernation premium, which is
$1.6\times$ the published value. Our Burst-HADS is therefore a more expensive
baseline under hibernation than the published one, and R-BurstHADS's cost
advantage over it is flattered by an amount this study has not bounded.

\textbf{The advantage depends on launch headroom.} Under limits R-BurstHADS is
the only scheduler with more than a single missed task, and its cost lead over
Burst-HADS falls from 14.6\,\% to 5.4\,\%.

\textbf{A retained addition favours the comparison.} The proactive-burstable
improvement guard is ours, not the published algorithm's, and removing it turns
R-BurstHADS's 5.3\,\% cost lead into a 5.8\,\% deficit.

\textbf{The sweep catalogue penalises the baseline's mechanism.} At current
prices a burst-mode rescue costs $1.15\times$ the on-demand alternative, against
$0.83\times$ in the published catalogue, so Burst-HADS's central mechanism is a
cost penalty in the main comparison.

\textbf{Per-type risk is never exercised.} The scenarios hibernate every spot
machine at the same rate, so the per-type hibernation differential that
Eq.~\eqref{eq:thm1} reasons about is never instantiated. No result in this paper
is attributed to risk awareness.

\textbf{Five deviations remain open}, one of them unmeasured.

\textbf{Scope.} One simulator, synthetic workloads, 30 seeds per cell.

%% ===================================================================
\section{Future work}
%% ===================================================================

Four items follow directly. The spot pool size is not stated
in~\cite{teylo2023} and its hibernation counts depend on it; a sensitivity study
over pool size would bound the validation gap. The provisioning and
replacement-utilisation diagnostics predate the freeze and should be re-run on
the frozen code. The one unmeasured deviation should be measured. And
heterogeneous per-type hibernation rates would, for the first time, exercise the
risk reasoning in Eq.~\eqref{eq:thm1}.

%% ===================================================================
\section{Conclusion}
%% ===================================================================

Burst-HADS buys makespan with cost. Across 80 cells under provider launch
limits it finishes 19.1\,\% sooner than HADS and pays 21.3\,\% more.
R-BurstHADS, which adds preemptive and deadline-sized reactive spot provisioning
around an otherwise unmodified Burst-HADS, finishes 27.9\,\% sooner than HADS at
a 12.9\,\% premium --- 12.7\,\% faster and 5.4\,\% cheaper than Burst-HADS, and
dominant on both axes in 44 of 75 cells. Without launch limits the margin widens
to 17.8\,\% and 14.6\,\%, in 69 of 80 cells.

Two results cut the other way and we report them as prominently as the first.
Under launch limits R-BurstHADS is the only scheduler with more than a single
missed task, 109 in 4.3\,\% of runs, and all of them disappear once the limits
are removed: the method depends on provisioning headroom that a provider may not
grant. And our re-implemented baselines do not reproduce the published
hibernation cost result, which we document, decompose and leave open rather than
resolve. The accompanying deviation register is offered as part of the
contribution: a comparison whose disagreements with its own sources are stated
is worth more than one whose are not.

\bibliographystyle{IEEEtran}
\begin{thebibliography}{9}

\bibitem{teylo2021}
L.~Teylo \emph{et al.}, ``A hibernation aware dynamic scheduler for cloud
environments,'' \emph{Cluster Computing}, 2021.

\bibitem{teylo2023}
L.~Teylo \emph{et al.}, ``Scheduling bag-of-tasks in clouds using spot and
burstable virtual machines,'' \emph{IEEE Transactions on Cloud Computing},
2023.

\end{thebibliography}

\appendices

\section{Deviation register}
The full register, with the measured effect of each entry, is reproduced from
\texttt{DEVIATIONS.md}.

\section{Per-cell results}
Per-cell means and 95\,\% confidence intervals for makespan and cost, all three
schedulers, both limit regimes; and the infeasible runs with their cause.

\section{Validation detail}
Comparison against the published Table~7 and Table~9, cell by cell; the
ruled-out causes of the cost gap; and the trajectory of the audit's
corrections.

\section{Reproducibility}
Frozen tag and code fingerprint, the results pack with per-file hashes and
commits, the generating script, and the pre-registered decision rule.

\end{document}
